% Chapter 9: Analytical Derivation
% ============================================================

\chapter{Analytical Derivation}
\label{ch:analytical}

This chapter presents the step-by-step algebraic derivation of the polynomial system underlying the VDPP. We parameterize the problem by the subtended angles at $P$, derive the resulting equations, and analyze their algebraic structure.

\section{The Parameterization}

Recall from Chapter~\ref{ch:formulation} that we place $P$ at the origin and parameterize by the subtended angles $\alpha = \angle BPC$, $\beta = \angle APC$, $\gamma = \angle APB$, with the constraint $\alpha + \beta + \gamma = 2\pi$.

Setting $\gamma = 2\pi - \alpha - \beta$, the side lengths are:
\begin{align}
a(\alpha) &= \sqrt{d_2^2 + d_3^2 - 2d_2 d_3 \cos\alpha}, \label{eq:a-of-alpha} \\
b(\beta) &= \sqrt{d_1^2 + d_3^2 - 2d_1 d_3 \cos\beta}, \label{eq:b-of-beta} \\
c(\alpha, \beta) &= \sqrt{d_1^2 + d_2^2 - 2d_1 d_2 \cos(2\pi - \alpha - \beta)}. \label{eq:c-of-alpha-beta}
\end{align}

Since $\cos(2\pi - x) = \cos(x)$:
\begin{equation}
c(\alpha, \beta) = \sqrt{d_1^2 + d_2^2 - 2d_1 d_2 \cos(\alpha + \beta)}.
\end{equation}

\section{Squaring to Eliminate Radicals}

Define the squared side lengths:
\begin{align}
A &\equiv a^2 = d_2^2 + d_3^2 - 2d_2 d_3 \cos\alpha, \label{eq:A-sq} \\
B &\equiv b^2 = d_1^2 + d_3^2 - 2d_1 d_3 \cos\beta, \label{eq:B-sq} \\
C &\equiv c^2 = d_1^2 + d_2^2 - 2d_1 d_2 \cos(\alpha + \beta). \label{eq:C-sq}
\end{align}

These are \emph{trigonometric polynomials} in $(\alpha, \beta)$.

\section{The Elimination Problem}

\subsection{From Angles to Side Lengths}

Given a specific $(\alpha, \beta) \in \Omega$, the computation of $(a, b, c)$ via~\eqref{eq:a-of-alpha}--\eqref{eq:c-of-alpha-beta} is straightforward. The \emph{inverse} problem---given $(a, b, c)$, find $(\alpha, \beta)$---requires solving:
\begin{align}
\cos\alpha &= \frac{d_2^2 + d_3^2 - a^2}{2d_2 d_3}, \label{eq:cos-alpha} \\
\cos\beta &= \frac{d_1^2 + d_3^2 - b^2}{2d_1 d_3}, \label{eq:cos-beta} \\
\cos(\alpha + \beta) &= \frac{d_1^2 + d_2^2 - c^2}{2d_1 d_2}. \label{eq:cos-gamma}
\end{align}

These three equations are not independent: the third is determined by the first two via the cosine addition formula:
\begin{equation}
\cos(\alpha + \beta) = \cos\alpha \cos\beta - \sin\alpha \sin\beta.
\end{equation}

\subsection{The Consistency Equation}

Substituting~\eqref{eq:cos-alpha} and~\eqref{eq:cos-beta} into the cosine addition formula and setting equal to~\eqref{eq:cos-gamma}:

\begin{equation}
\label{eq:consistency}
\frac{d_1^2 + d_2^2 - c^2}{2d_1 d_2} = \frac{(d_2^2 + d_3^2 - a^2)(d_1^2 + d_3^2 - b^2)}{4d_1 d_2 d_3^2} - \frac{\sqrt{(4d_2^2 d_3^2 - (d_2^2 + d_3^2 - a^2)^2)(4d_1^2 d_3^2 - (d_1^2 + d_3^2 - b^2)^2)}}{4d_1 d_2 d_3^2}.
\end{equation}

This is the \emph{consistency equation}: it must hold for any valid triple $(a, b, c)$ arising from the VDPP with distances $(d_1, d_2, d_3)$.

\subsection{Eliminating the Square Root}

To obtain a polynomial relation, isolate the square root term and square both sides:

Let:
\begin{align}
u &= d_2^2 + d_3^2 - a^2, \label{eq:u-def} \\
v &= d_1^2 + d_3^2 - b^2, \label{eq:v-def} \\
w &= d_1^2 + d_2^2 - c^2. \label{eq:w-def}
\end{align}

Then~\eqref{eq:consistency} becomes:
\begin{equation}
2d_3^2 w = uv - \sqrt{(4d_2^2 d_3^2 - u^2)(4d_1^2 d_3^2 - v^2)}.
\end{equation}

Rearranging and squaring:
\begin{equation}
(uv - 2d_3^2 w)^2 = (4d_2^2 d_3^2 - u^2)(4d_1^2 d_3^2 - v^2).
\end{equation}

\begin{theorem}[The Master Polynomial]
\label{thm:master-polynomial}
Expanding and simplifying, the consistency equation yields:
\begin{equation}
\label{eq:master-poly}
u^2 v^2 - 4d_3^2 uvw + 4d_3^4 w^2 = 4d_2^2 d_3^2 \cdot 4d_1^2 d_3^2 - 4d_2^2 d_3^2 v^2 - 4d_1^2 d_3^2 u^2 + u^2 v^2.
\end{equation}

After cancellation:
\begin{equation}
\label{eq:master-simplified}
- 4d_3^2 uvw + 4d_3^4 w^2 = 16d_1^2 d_2^2 d_3^4 - 4d_2^2 d_3^2 v^2 - 4d_1^2 d_3^2 u^2.
\end{equation}

Dividing by $4d_3^2$:
\begin{equation}
\label{eq:master-final}
\boxed{d_3^2 w^2 - uvw + d_2^2 v^2 + d_1^2 u^2 - 4d_1^2 d_2^2 d_3^2 = 0,}
\end{equation}
where $u$, $v$, $w$ are defined by~\eqref{eq:u-def}--\eqref{eq:w-def}.
\end{theorem}

\begin{remark}
Equation~\eqref{eq:master-final} is equivalent to the vanishing of the Cayley--Menger determinant $\CM_4 = 0$, confirming that both approaches yield the same polynomial constraint.
\end{remark}

\section{Degree Analysis}

Substituting $u = d_2^2 + d_3^2 - a^2$, $v = d_1^2 + d_3^2 - b^2$, $w = d_1^2 + d_2^2 - c^2$ into~\eqref{eq:master-final}:

\begin{itemize}
    \item Each of $u, v, w$ is a polynomial of degree 2 in the side lengths $a, b, c$ (after treating $d_1, d_2, d_3$ as constants).
    \item The terms $u^2$, $v^2$, $w^2$ are degree 4 in side lengths.
    \item The term $uvw$ is degree 6 in side lengths.
\end{itemize}

Therefore, the master polynomial~\eqref{eq:master-final} is of \textbf{degree 6} in $(a, b, c)$, or equivalently, degree 3 in $(a^2, b^2, c^2)$.

\section{The Sextic Polynomial}

For fixed $(d_1, d_2, d_3)$ and with one additional constraint (e.g., fixing one side length), the master polynomial becomes a univariate polynomial. The degree of this polynomial depends on the constraint:

\begin{itemize}
    \item \textbf{Equilateral constraint} ($a = b = c = s$): Degree 6 in $s$, but reduces to degree 3 in $s^2$, hence a cubic in $s^2$.
    \item \textbf{Isosceles constraint} ($b = c$): Degree 6 in $a$ (with $b$ related to $a$ through the constraint).
    \item \textbf{One side fixed}: Degree 6 in the remaining two sides.
\end{itemize}

\section{Perimeter Bounds from the Analytical Framework}

From the parameterization~\eqref{eq:a-of-alpha}--\eqref{eq:c-of-alpha-beta}, we can derive bounds on each side length:

\begin{proposition}[Side Length Bounds]
For each side, the extreme values over the feasible domain are:
\begin{align}
\abs{d_2 - d_3} &< a < d_2 + d_3, \\
\abs{d_1 - d_3} &< b < d_1 + d_3, \\
\abs{d_1 - d_2} &< c < d_1 + d_2.
\end{align}
The upper bounds are achieved as the corresponding angle at $P$ approaches $\pi$; the lower bounds are achieved as it approaches $0$.
\end{proposition}

\begin{corollary}[Perimeter Bounds]
\begin{equation}
\perim_{\min} \leq \perim \leq \perim_{\max},
\end{equation}
where:
\begin{align}
\perim_{\max} &< (d_2 + d_3) + (d_1 + d_3) + (d_1 + d_2) = 2(d_1 + d_2 + d_3), \\
\perim_{\min} &> \abs{d_2 - d_3} + \abs{d_1 - d_3} + \abs{d_1 - d_2}.
\end{align}
The exact extremal values are attained at boundary configurations of the feasible domain $\Omega$.
\end{corollary}
